\cotefile{57}
\scanpage{87}
ouvert \deleted{\unsure{de|un}} $U$ de $X$, la donnée d'une connexion sur $\underline L=\underline O_X(D)$, i.e d'un splittage de
\[
0\mathbin{-}\underline L\otimes\underline\Omega^1_{X/S}\to P^1(\underline L)\mathbin{-}\underline L\mathbin{-}0,
\]
i.e. d'une projection $P^1(\underline L)\mathbin{-}\underline L\otimes\underline\Omega^1_{X/S}$, ou enfin d'un op. 1-différentiel rel$/S$
$\underline L\xrightarrow{\unsure{d'_L|d_L}}\underline L\otimes\Omega^1_{X/S}$ satisfaisant $d_L(\lambda x)=\deleted{\unsure{d\lambda}}x\otimes d(\lambda)+\lambda d_L(x)$,
\margin{\deleted{\unsure{dé}}}\deleted{\unsure{équiva|équivaut}} définit, grâce à la section rationnelle relative canonique $\varphi_{\unsure{D}}$ de $\underline L=\underline O_X(D)$, une \unsure{élémt|élément} $\varpi$
\[
\frac{d_L(\varphi_D)}{\varphi_D}\in\underline R_{\unsure{U/S|X/S}}\bigl(\underline L^{-1}\otimes\underline L\otimes\Omega^1_{X/S}\bigr)=\underline R_{U/S}\bigl(\underline\Omega^1_{X/S}\bigr),
\]
dont la connaissance détermine $d_L$ donc la connexion, car on aura \inserted{sur $U-\operatorname{supp}D$}
\[
d_L(\lambda\varphi_{\unsure{D}})=\varphi\otimes d(\lambda)+\lambda\underbrace{d_L(\varphi)}=\varphi\otimes d\lambda+\lambda\varphi\varpi.
\]

\deleted{\unsure{Donc}} D'autre part, quand une $\varpi$ provient-elle d'une connexion ? On voit tout de suite que c'est sss
\[
\varpi\deleted{\unsure{=}}-\frac{df}{f}\in\inserted{\Gamma(\underline\Omega^1_{V/S})}\ \deleted{\text{\unsure{régulière}}},
\]
pour \deleted{\unsure{toute}} \deleted{\unsure{section|intégrale}}\inserted{équation} $f$ de $D$ sur un ouvert $V\subset U$. Donc on trouve que $\underline E_{X/S}$ s'interprète comme le faisceau des $(D,c)$, $D$ un diviseur, $c$ une connexion sur $\underline O_X(D)$. Le morphisme $\alpha$ est maintenant évident. Si $f_*(\underline O_X^{*})\xrightarrow{\varphi\mapsto\frac{d\varphi}{\varphi}}f_*(\underline\Omega^1_{X/S})$ est nul, i.e. $\underline\omega_{X/S}\mathbin{-}\underline\omega_{X/S}/\operatorname{Im}f_*(\underline O_X^{*})$ un isom, on peut donc interpréter l'ext. $f_*(\underline E_{X/S})$ comme image inverse de l'ext $\underline{\operatorname{Pic}}_{X/S}^{w}$.
% archive sheet number: 86
